% Adapted from the defense at 3593fbf25434c4715b37537eae1287bf49239fa7; original blob 9cca660433532a16af3cf736cf9383184d3ba9e6.
                \begin{tikzpicture}[xscale=1.0, yscale=1.0]
                    \useasboundingbox (-0.45, -2.85) rectangle (13.35, 2.25);

                    \tikzset{pt/.style={circle, fill=inria-2024-bleu-canard, inner sep=0pt, minimum size=2.2pt}}
                    \tikzset{surf/.style={inria-rouge, line width=1.6pt, line join=round, line cap=round}}
                    \tikzset{carrosserie/.style={draw=gray!45, line width=0.7pt, fill=gray!8, line join=round}}
                    \tikzset{rayon/.style={gris_fonce_inria, line width=0.35pt, opacity=0.40}}
                    \tikzset{cadre/.style={draw=gray!30, line width=0.9pt, rounded corners=3pt}}

                    \newcommand{\voiture}{%
                        \draw[carrosserie]
                            (0.00,0.00) -- (0.00,0.42) -- (0.48,0.46) -- (0.92,0.98)
                            -- (2.00,0.98) -- (2.42,0.46) -- (3.00,0.42) -- (3.00,0.00) -- cycle;
                        \draw[gray!45, line width=0.7pt] (0.70,0.00) circle (0.20);
                        \draw[gray!45, line width=0.7pt] (2.30,0.00) circle (0.20);
                    }
                    \newcommand{\contour}{(0.00,0.42) (0.48,0.46) (0.92,0.98) (2.00,0.98) (2.42,0.46) (3.00,0.42)}

                    % ================== Panneau gauche : de près ==================
                    \begin{scope}[xshift=0.4cm, yshift=0.35cm]
                        \node[cadre, minimum width=5.0cm, minimum height=2.30cm] at (1.5,0.60) {};
                        \voiture
                        % Nappes de balayage tous les 0.108 : l'une longe le capot et le coffre
                        % (0.44), une autre le toit (0.98) ; chaque retour est à l'intersection
                        % d'une nappe et du contour (pas 0.108 le long d'une nappe).
                        \foreach \y in {0.224, 0.332, 0.440, 0.548, 0.656, 0.764, 0.872, 0.980, 1.088} {\draw[rayon] (-0.55,\y) -- (3.55,\y);}
                        \foreach \p in {(0.020,0.440),(0.128,0.440),(0.236,0.440),(0.344,0.440),(0.452,0.440),(0.555,0.548),(0.646,0.656),(0.738,0.764),(0.829,0.872),(0.980,0.980),(1.088,0.980),(1.196,0.980),(1.304,0.980),(1.412,0.980),(1.520,0.980),(1.628,0.980),(1.736,0.980),(1.844,0.980),(1.952,0.980),(2.087,0.872),(2.174,0.764),(2.261,0.656),(2.348,0.548),(2.450,0.440),(2.558,0.440),(2.666,0.440),(2.774,0.440),(2.882,0.440),(2.990,0.440)} {
                            \node[pt] at \p {};
                        }
                        \node[font=\fontsize{10}{11}\selectfont, text=gris_fonce_inria, anchor=north] at (1.5,-0.44) {\textbf{Near} object};
                    \end{scope}

                    % ================== Panneau droit : de loin ==================
                    \begin{scope}[xshift=6.9cm, yshift=0.35cm]
                        \node[cadre, minimum width=5.0cm, minimum height=2.30cm] at (1.5,0.60) {};
                        \voiture
                        % Nappes 2,5 fois plus espacées (0.27) et retours 3 fois plus espacés
                        % le long d'une nappe (0.34) : neuf retours sur trois nappes.
                        \foreach \y in {0.17, 0.44, 0.71, 0.98, 1.25} {\draw[rayon] (-0.55,\y) -- (3.55,\y);}
                        \foreach \p in {(0.080,0.440),(0.420,0.440),(0.692,0.710),(1.100,0.980),(1.440,0.980),(1.780,0.980),(2.218,0.710),(2.580,0.440),(2.920,0.440)} {
                            \node[pt] at \p {};
                        }
                        \node[font=\fontsize{10}{11}\selectfont, text=gris_fonce_inria, anchor=north] at (1.5,-0.44) {\textbf{Far} object};
                    \end{scope}

                    % ================== Flèches ==================
                    \draw[-{Latex[length=2.8mm]}, line width=1.6pt, color=gray!45] (1.90,-0.95) -- (1.90,-1.28);
                    \draw[-{Latex[length=2.8mm]}, line width=1.6pt, color=gray!45] (8.40,-0.95) -- (8.40,-1.28);

                    % ================== Surfaces reconstruites ==================
                    % à gauche : fine ; à droite : grossière, mais même cible
                    \begin{scope}[xshift=0.4cm, yshift=-2.40cm]
                        \draw[gray!45, line width=0.7pt, densely dotted] plot coordinates {\contour};
                        \draw[surf] plot coordinates {\contour};
                    \end{scope}
                    \begin{scope}[xshift=6.9cm, yshift=-2.40cm]
                        \draw[gray!45, line width=0.7pt, densely dotted] plot coordinates {\contour};
                        % La surface interpole les neuf retours observés : ses sommets sont
                        % sur le contour gris, les cordes coupent les angles. Grossière, mais
                        % c'est bien le même objet, c'est tout le propos de la planche.
                        \draw[surf] plot coordinates {(0.080,0.440) (0.420,0.440) (0.692,0.710) (1.100,0.980) (1.440,0.980) (1.780,0.980) (2.218,0.710) (2.580,0.440) (2.920,0.440)};
                    \end{scope}

                    % Both comparison signs share the midpoint between the two panels.
\def\comparisonx{5.15}
\node[font=\fontsize{18}{20}\selectfont\bfseries, text=inriablue] at (\comparisonx,0.95) {$\ne$};
\node[font=\fontsize{18}{20}\selectfont\bfseries, text=inria-rouge] at (\comparisonx,-1.86) {$\approx$};

                    \node[font=\fontsize{8}{9}\selectfont, text=inria-2024-bleu-canard, anchor=east, align=right] at (12.45,0.95)
                        {\textbf{Sampling}\\depends on\\sensor and range};
                    \node[font=\fontsize{8}{9}\selectfont, text=inria-rouge, anchor=east, align=right] at (12.45,-1.90)
                        {\textbf{Geometry}\\may vary less\\(hypothesis)};
                \end{tikzpicture}
